% from ega4-2-fr.tex
\oldpage[IV-2]{217}

excellent, donc vérifie la condition (6.12.4, $a)$) en vertu de (7.8.3, (v));
cela prouve que $A$ lui-même vérifie la condition de (6.12.4, $a)$), donc aussi
(6.12.4, $c)$), qui n'est autre que la condition (7.8.2, (iii)). Pourtant $A$
ne satisfait pas à la condition (7.8.2, (i)), puisque nous avons vu (5.6.11)
qu'il n'est pas universellement caténaire.

(iii) On peut remplacer la condition (i) de la définition (7.8.2) par la
suivante (plus faible en apparence, compte tenu de (5.6.10)) :

(i$'$) \emph{Pour tout idéal premier minimal} $\mathfrak p_i$ \emph{de} $A$
$(1\leq i\leq r)$, \emph{soit} $A'_i$ \emph{la clôture intégrale de l'anneau
intègre} $A_i=A/\mathfrak p_i$; \emph{alors, pour tout idéal maximal}
$\mathfrak m'$ \emph{de} $A'_i$, \emph{on a, en désignant par} $\mathfrak m$
\emph{l'image réciproque de} $\mathfrak m'$ \emph{dans} $A$,
\begin{env}[7.8.4.1]
\label{IV.2.s7.8.4.1-fr}
\[
\dim(A_{\mathfrak m})=\dim((A'_i)_{\mathfrak m'})
\]
\end{env}

En effet, on a vu (Remarque (i)) que sous les seules hypothèses (ii) et (iii)
de (7.8.2), les \emph{analogues} des propriétés (ii), (v), (vi) et (vii) de
(7.8.3) sont valables. On déduit donc d'abord de (vi) et (ii) que les $A'_i$
vérifient (7.8.2, (ii) et (iii)); il résulte alors de (v) que les complétés
$((A'_i)_{\mathfrak m'})^\wedge$ sont intègres (et intégralement clos). Soit
maintenant $\mathfrak m$ un idéal maximal quelconque de $A$; pour tout $i$ tel
que $\mathfrak p_i\subset\mathfrak m$, la clôture intégrale de
$A_{\mathfrak m}/\mathfrak p_iA_{\mathfrak m}$ est une
$(A_{\mathfrak m}/\mathfrak p_iA_{\mathfrak m})$-algèbre finie par (vi), donc
semi-locale, et ses composants locaux sont de la forme
$(A'_i)_{\mathfrak m'}$, où $\mathfrak m'$ est un idéal premier
(nécessairement maximal) de $A'_i$ au-dessus de
$\mathfrak m/\mathfrak p_i$; il résulte alors de (vii), de ce qui précède et de
(7.8.4.1) que les quotients du complété
$(A_{\mathfrak m}/\mathfrak p_iA_{\mathfrak m})^\wedge$ par ses idéaux
premiers minimaux ont tous même dimension égale à $\dim(A_{\mathfrak m})$. Par
suite (7.1.9 et 7.1.8, $b)$), $A_{\mathfrak m}$ est formellement caténaire, et
\emph{a fortiori} (7.1.11) universellement caténaire; il en est donc de même de
$A_{\mathfrak p}$ pour tout idéal premier $\mathfrak p$ de $A$ (5.6.3, (i)),
donc aussi de $A$.

En particulier, si $A$ est normal, ou plus généralement si $A_{\mathfrak m}$
est unibranche pour tout idéal maximal $\mathfrak m$ de $A$, la condition
(i$'$) est toujours remplie, puisque $\mathfrak m$ ne peut contenir qu'un seul
idéal premier $\mathfrak p_i$ (0, 16.1.5); on peut dans ce cas omettre (i) dans
la définition (7.8.2). Plus particulièrement, si $A$ est un anneau noethérien
\emph{local unibranche}, il est excellent si et seulement si ses fibres
formelles sont géométriquement régulières.

(iv) Rappelons qu'un anneau quotient d'un anneau local de Cohen-Macaulay (et
\emph{a fortiori} un anneau quotient d'un anneau local régulier) vérifie aussi
les propriétés (7.8.3, (ix) et (x)) en vertu de (7.2.7); mais l'intérêt des
quotients d'anneaux réguliers réside surtout dans leurs propriétés
cohomologiques (chap. III, 3\textsuperscript{e} Partie).

(v) Nous verrons plus loin (\S~18) que si $k$ est un corps valué complet non
discret, l'anneau $k\{\!\{T_1,\ldots,T_n\}\!\}$ des \emph{séries entières
convergentes} (dans un voisinage de $0$ dans $k^n$) est un \emph{anneau
excellent}.
\end{remarks}

\begin{env}[7.8.5]
\label{IV.2.s7.8.5-fr}
On dit qu'un préschéma localement noethérien $X$ est \emph{excellent} si pour
un recouvrement $(U_\alpha)$ de $X$ formé d'ouverts affines, chacun des anneaux
$A_\alpha$ des $U_\alpha$ est excellent; cette propriété est alors vraie pour
\emph{tout} recouvrement $(U_\alpha)$ de $X$ formé d'ouverts affines.
\end{env}

\begin{proposition}[7.8.6]
\EGAReaderTarget{ega4.c04.prop.7-8-6}{IV.2.s7.8.6-fr}
Soit $X$ un préschéma excellent.
\begin{enumerate}
\item[\emph{(i)}] Si $f:X'\to X$ est un morphisme localement de type fini, $X'$ est
excellent.

